%\setcounter{page}{67}
%\fancyhead[ER]{\emph{VIII. Beweis der Chudnovsky-Formel}}



% Neu in v2: Übersetzung, jetzt auch englisch

% Neu in v3: Anhänge, Kap 10 umstrukturiert.

% Neu in v4: dreimal 1944000 statt dreimal 1994000 im Beweis von 5.14
% Neu in v4: introduction "at least 14 digits" durch "on average 14 digits" ersetzt.

% Neu in v5: Introduction 1728 vor J(tau_163) ergänzt
% Neu in v5: Introduction massiv umformuliert, Geschichtsquellen [1] und [21] und die Grafik ergänzt

% Neu in v6:
% - Quelle bei den Chudnovskys ist Gleichung 1.4 und nicht 4.1 !
%   (Thm. 0.1 und 9.7 waren richtig, thm 10.11 war falsch); NUR die Formel mit den Koeffizienten ist 1.5.
% - Lückenbreite in Anhängen A und B (z.B. Bemerkung A.6) optimiert (Leerzeichen um \mathbb Z\left[...\right] kleiner/größer)
% - Verlängerung der Linie im Diagramm der Einführung (auf aktuelles Datum anpassen!)
% - defltilde in deutscher Version Grammatik geändert (Def. 8.3)
% - Werte von \tau_N für gerade N waren falsch in Thm. 10.1 und Tabelle 10.1
% - Def. 3.4 "equivalent under modular transformations" ergänzt
% - Näherungen-Seite rausgeworfen
%   - Beweis von 9.7 entsprechend leicht angepasst
%   - Bemerkung 5.9 entsprechend angepasst
%   - Werte in Kap. 10 ergänzt
% - Anhang B: "ganzalgebraisch" statt "ganzalgebraisch in Z" geschrieben (mehrfach)
% - Kapitel 10 komplett neu
%   - Abb. 8.1: Bildunterschrift entsprechend angepasst
%   - Fußnote in Einleitung entsprechend angepasst
% - Hinweis auf Thm. 10.13 (Dezimalen) in Einleitung ergänzt



% - Man könnte noch in Bemerkung 9.3 ergänzen wie man das beweist
%     (Legendre-Relation 2.4 und 3.8 mit a=(-?)1/tau und 4.5 (eta_1 ~ E_2))
